%===============================================================================
% Kapitel 1: Einleitung
%===============================================================================

\chapter{Einleitung}
\label{chap:introduction}

Die Einleitung führt den Leser in das Thema ein und motiviert die Arbeit.
Sie sollte folgende Punkte abdecken:

\begin{enumerate}
    \item \textbf{Kontext}: Allgemeiner Hintergrund des Themas
    \item \textbf{Problem}: Was ist das konkrete Problem?
    \item \textbf{Lücke}: Welche Forschungslücke besteht?
    \item \textbf{Zielsetzung}: Was wird in dieser Arbeit untersucht?
    \item \textbf{Struktur}: Wie ist die Arbeit aufgebaut?
\end{enumerate}

\section{Hintergrund}
\label{sec:intro:background}

Erläutere hier den wissenschaftlichen Hintergrund und den Stand der Forschung.
Zitiere relevante Arbeiten anderer Autoren.

\section{Problemstellung}
\label{sec:intro:problem}

Formuliere das zu lösende Problem klar und präzise.

\section{Ziele und Forschungsfragen}
\label{sec:intro:objectives}

Definiere die genauen Ziele dieser Arbeit. Waren zutreffend:
\begin{itemize}
    \item Forschungsfrage 1
    \item Forschungsfrage 2
    \item Forschungsfrage 3
\end{itemize}

\section{Struktur der Arbeit}
\label{sec:intro:structure}

Die vorliegende Arbeit ist wie folgt strukturiert:
\begin{itemize}
    \item \chref{chap:methodology} beschreibt die Methodik
    \item \chref{chap:results} präsentiert die Ergebnisse
    \item \chref{chap:discussion} diskutiert die Befunde
    \item \chref{chap:conclusion} zieht abschließende Schlussfolgerungen
\end{itemize}
